% ---------------------------------------------------------------------------
% Supplementary Information: land-surface validation at eddy-covariance towers.
%
% Drop-in section for the legoESM supplement.  Uses only packages the supplement
% already loads (graphicx, subcaption, booktabs, natbib, amsmath).  The figure
% expects ec_site_combined.png in the manuscript figure directory; it is checked
% in at docs/land/figures/ec_site/ec_site_combined.png.
%
% Regenerate the figure with
%   python scripts/plot/plot_ec_site_publication.py \
%       scripts/validate/ec_site_example_data <out_prefix>
% See docs/land/ec_site_evaluation_runbook.md.
%
% BibTeX entries for every key cited below are collected at the end of this file.
% Bibliographic details were verified against the publisher records; the schemes
% themselves were verified against the implementation.
% ---------------------------------------------------------------------------

\section{Land-surface validation at eddy-covariance flux towers}
\label{si:ecsite}

The land component is evaluated offline against point-scale eddy-covariance
measurements, which constrain the coupled water, energy and carbon fluxes at the
sub-daily to seasonal timescales that a coupled simulation cannot isolate. This
section documents the sites, the forcing data, the model configuration, and the
experimental and evaluation protocols.

\subsection{Sites}

Four FLUXNET towers were selected to span the temperate and Mediterranean
regimes over which the two-big-leaf canopy is well posed
(Table~\ref{tab:ecsites}): two temperate deciduous forests differing in climate
and soil (US-MMS, DE-Hai), a montane evergreen needleleaf forest (DE-Obe), and a
water-limited oak savanna (US-Ton). Peak leaf area index ranges from
2.9\,m$^2$\,m$^{-2}$ at the savanna to 6.2\,m$^2$\,m$^{-2}$ at the closed
forests, and soil texture spans clay to loam, so the four sites exercise
contrasting water-limitation, phenology and surface-roughness regimes.

\begin{table}[h]\centering\footnotesize
\caption{Eddy-covariance sites used for the land-surface validation. Canopy
height and tower height are site metadata; the tower height sets the reference
level of the surface-layer profile in the model.}\label{tab:ecsites}
\begin{tabular}{@{}llrrrrrl@{}}
\toprule
Site & Biome (PFT) & Lat. & Lon. & $h_c$ & $z_{\mathrm{ref}}$ & LAI$_{\max}$ & Period \\
     &             & ($^\circ$) & ($^\circ$) & (m) & (m) & (m$^2$\,m$^{-2}$) & \\
\midrule
US-MMS & temperate deciduous forest (DBF) &  39.32 & $-86.41$ & 27 & 46.0 & 6.2 & 2015--2018 \\
DE-Obe & montane needleleaf forest (ENF)  &  50.79 &   13.72  & 20 & 30.0 & 5.2 & 2015--2018 \\
US-Ton & Mediterranean oak savanna (SAV)  &  38.43 & $-120.97$ & 10 & 23.5 & 2.9 & 2015--2018 \\
DE-Hai & humid deciduous forest (DBF)     &  51.08 &   10.45  & 34 & 42.0 & 6.2 & 2010--2013 \\
\bottomrule
\end{tabular}
\end{table}

\subsection{Forcing data}

Each site is driven by a per-site forcing file combining FLUXNET2015 FULLSET
meteorology\cite{Pastorello2020} (downwelling shortwave and longwave radiation,
air temperature, vapour-pressure deficit, wind speed, surface pressure,
precipitation and CO$_2$ mole fraction) with MODIS leaf area
index\cite{Myneni2002} and albedo\cite{Schaaf2002}, the diffuse fraction of
photosynthetically active radiation from the BESS radiation
product\cite{Ryu2018}, and site parameters. The forcing is half-hourly except at
US-MMS, which is hourly. Gaps in the meteorological forcing are filled by FLUXNET
\texttt{\_F\_MDS} substitution, with gaps of at least 24\,h filled from a
month-by-time-of-day climatology and shorter interior gaps interpolated linearly.
Gap-filling is applied only to the forcing: the observed fluxes used as
evaluation targets are never gap-filled and retain their gaps, so no filled value
can enter a skill statistic.

\subsection{Model configuration}

The land column is run in its multi-layer configuration. Soil moisture follows
the mass-conservative mixed-form Richards equation\cite{Celia1990} with
van~Genuchten retention\cite{VanGenuchten1980} whose parameters are derived from
the site soil texture, solved by Picard iteration on a tridiagonal system. Soil
temperature is advanced by backward-Euler heat diffusion with Johansen thermal
conductivity\cite{Johansen1975}. The column is discretised into eight
geometrically spaced layers reaching approximately 6.4\,m (extended to 10\,m at
the deep-rooted savanna), with free drainage at the base.

Surface exchange uses the two-big-leaf canopy scheme. Shortwave transfer follows
a two-stream partition into sunlit and shaded leaves\cite{Sellers1985,Ryu2011},
each of which carries Farquhar--von~Caemmerer--Berry C$_3$
photosynthesis\cite{Farquhar1980} with Kattge--Knorr thermal
acclimation\cite{KattgeKnorr2007}, or C$_4$ photosynthesis following
Collatz\cite{Collatz1992}, coupled to Ball--Berry stomatal
conductance\cite{Ball1987}. The resulting six-variable system in sunlit and
shaded leaf temperature, the two intercellular CO$_2$ concentrations, and
canopy-air temperature and humidity is closed by Newton--Raphson iteration; an
outer Picard loop reconciles the canopy solution with the soil thermal solver.
Aerodynamic exchange is referenced to the measured tower height at each site,
with roughness length and displacement height taken from a plant-functional-type
table. The Ball--Berry cuticular intercept is not down-regulated by soil-moisture
stress, so a live canopy retains a baseline conductance while a dormant canopy
self-limits as its leaf area declines.

Rooting depth and soil-column depth are prescribed modelling choices rather than
measurements: rooting depth is the e-folding depth of an exponential
root-density profile, set between 1 and 5\,m according to the plant functional
type and soil profile, deepest at the phreatophytic savanna.

\subsection{Experimental protocol}

Each site is integrated as an independent offline single column at the native
forcing timestep over the full analysis window, with soil moisture and soil
temperature carried prognostically. The soil state is initialised from the first
observed soil temperature and volumetric water content; no separate spin-up cycle
is performed, so the earliest part of each record reflects adjustment away from
the observed initial state.

\subsection{Evaluation protocol}

A timestep enters the evaluation only if every required forcing field was
genuinely observed before gap-filling, which retains 63, 83, 96 and 60\,\% of
timesteps at US-MMS, DE-Obe, US-Ton and DE-Hai respectively. Model and
observations are additionally co-sampled, that is, compared only on timesteps at
which both are available, so that a model--observation difference cannot arise
from the two being averaged over different samples. Sub-daily values are
aggregated to daily means wherever at least 30\,\% of a day is present, and skill
is summarised by the Nash--Sutcliffe efficiency, the coefficient of determination
and the mean bias.

Raw eddy-covariance fluxes under-close the surface energy budget by 22--34\,\% at
these sites, whereas the model closes it by construction. Skill is therefore
reported against both the raw and the energy-balance-closure-corrected
observations, which bracket the true flux. Pooled across sites, the modelled
turbulent fluxes are biased $+17$\,\% against the raw observations and $-17$\,\%
against the corrected observations, and so lie near the centre of the
closure-uncertainty envelope; the remaining bias in gross primary productivity,
which carries no closure ambiguity, is $+6$\,\%. Interannual variability is not
assessed: with two to four years per site, year-to-year means are dominated by
sampling noise.

\begin{figure}[h]\centering
\includegraphics[width=\textwidth]{ec_site_combined.png}
\caption{\textbf{Land-surface validation against eddy-covariance flux towers.}
The offline land model with the two-big-leaf canopy scheme at four FLUXNET sites:
US-MMS (mesic deciduous forest), DE-Obe (montane needleleaf forest), US-Ton (oak
savanna) and DE-Hai (humid deciduous forest). One row per site (a--t): the mean
summer diurnal cycle and the seasonal cycle of latent (orange) and sensible (blue)
heat; the GPP seasonal cycle; the modelled latent-heat partition into
transpiration (green) and soil evaporation (orange) with the observed total; and
soil moisture at the observed sensor depth and two deeper model layers. Lines are
modelled, markers are raw eddy-covariance observations, and shaded bands span up
to the energy-balance-closure-corrected value; model and observations are compared
only on matched timesteps. Pooled daily Nash--Sutcliffe efficiency against raw
observations is 0.58 (latent heat), 0.57 (sensible heat) and 0.80 (GPP), and 0.68
for both energy fluxes against the closure-corrected observations. Raw fluxes
under-close the energy budget by 22--34\,\% here, so the modelled turbulent fluxes
sit near the centre of the closure-uncertainty envelope.}\label{si:ecsite_fig}
\end{figure}

% ===========================================================================
% BibTeX entries for the keys cited above.  Add any that the main .bib lacks;
% VanGenuchten1980 and Farquhar1980 are already used elsewhere in the supplement.
% ===========================================================================
%
% @article{Pastorello2020,
%   author  = {Pastorello, Gilberto and Trotta, Carlo and Canfora, Eleonora and others},
%   title   = {The {FLUXNET2015} dataset and the {ONEFlux} processing pipeline for
%              eddy covariance data},
%   journal = {Scientific Data}, volume = {7}, pages = {225}, year = {2020},
%   doi     = {10.1038/s41597-020-0534-3}}
%
% @article{Myneni2002,
%   author  = {Myneni, R. B. and Hoffman, S. and Knyazikhin, Y. and others},
%   title   = {Global products of vegetation leaf area and fraction absorbed
%              {PAR} from year one of {MODIS} data},
%   journal = {Remote Sensing of Environment}, volume = {83}, pages = {214--231},
%   year    = {2002}, doi = {10.1016/S0034-4257(02)00074-3}}
%
% @article{Schaaf2002,
%   author  = {Schaaf, Crystal B. and Gao, Feng and Strahler, Alan H. and others},
%   title   = {First operational {BRDF}, albedo nadir reflectance products from
%              {MODIS}},
%   journal = {Remote Sensing of Environment}, volume = {83}, pages = {135--148},
%   year    = {2002}, doi = {10.1016/S0034-4257(02)00091-3}}
%
% @article{Ryu2018,
%   author  = {Ryu, Youngryel and Jiang, Chongya and Kobayashi, Hideki and Detto, Matteo},
%   title   = {{MODIS}-derived global land products of shortwave radiation and
%              diffuse and total photosynthetically active radiation at 5 km
%              resolution from 2000},
%   journal = {Remote Sensing of Environment}, volume = {204}, pages = {812--825},
%   year    = {2018}, doi = {10.1016/j.rse.2017.09.021}}
%   % Product: BESS Rad v1, https://developers.google.com/earth-engine/datasets/
%   %          catalog/SNU_ESL_BESS_Rad_v1
%
% @article{Ryu2011,
%   author  = {Ryu, Youngryel and Baldocchi, Dennis D. and Kobayashi, Hideki and others},
%   title   = {Integration of {MODIS} land and atmosphere products with a
%              coupled-process model to estimate gross primary productivity and
%              evapotranspiration from 1 km to global scales},
%   journal = {Global Biogeochemical Cycles}, volume = {25}, pages = {GB4017},
%   year    = {2011}, doi = {10.1029/2011GB004053}}
%
% @article{Celia1990,
%   author  = {Celia, Michael A. and Bouloutas, Efthimios T. and Zarba, Rebecca L.},
%   title   = {A general mass-conservative numerical solution for the unsaturated
%              flow equation},
%   journal = {Water Resources Research}, volume = {26}, pages = {1483--1496},
%   year    = {1990}, doi = {10.1029/WR026i007p01483}}
%
% @article{Sellers1985,
%   author  = {Sellers, P. J.},
%   title   = {Canopy reflectance, photosynthesis and transpiration},
%   journal = {International Journal of Remote Sensing}, volume = {6},
%   pages   = {1335--1372}, year = {1985}, doi = {10.1080/01431168508948283}}
%
% @article{KattgeKnorr2007,
%   author  = {Kattge, Jens and Knorr, Wolfgang},
%   title   = {Temperature acclimation in a biochemical model of photosynthesis:
%              a reanalysis of data from 36 species},
%   journal = {Plant, Cell \& Environment}, volume = {30}, pages = {1176--1190},
%   year    = {2007}, doi = {10.1111/j.1365-3040.2007.01690.x}}
%
% @article{Collatz1992,
%   author  = {Collatz, G. James and Ribas-Carbo, Miquel and Berry, Joseph A.},
%   title   = {Coupled photosynthesis-stomatal conductance model for leaves of
%              {C4} plants},
%   journal = {Australian Journal of Plant Physiology}, volume = {19},
%   pages   = {519--538}, year = {1992}, doi = {10.1071/PP9920519}}
%
% @incollection{Ball1987,
%   author    = {Ball, J. Timothy and Woodrow, Ian E. and Berry, Joseph A.},
%   title     = {A model predicting stomatal conductance and its contribution to
%                the control of photosynthesis under different environmental
%                conditions},
%   booktitle = {Progress in Photosynthesis Research}, editor = {Biggins, J.},
%   volume    = {4}, pages = {221--224}, publisher = {Martinus Nijhoff},
%   year      = {1987}, doi = {10.1007/978-94-017-0519-6_48}}
%
% @phdthesis{Johansen1975,
%   author  = {Johansen, Oistein},
%   title   = {Thermal conductivity of soils},
%   school  = {Norwegian University of Science and Technology, Trondheim},
%   year    = {1975},
%   note    = {English translation: CRREL Draft Translation 637, 1977}}
% ===========================================================================
